%*********************************************************************
\chapter{Conclusions}
\label{ch:conclusions}
%*********************************************************************

A stage-top incubator was designed, modelled, built and tested on a single STM32 microcontroller. The chamber air temperature held $37.028 \pm 0.038$\,\textdegree{}C over 80 minutes and 39\,\textdegree{}C to within 0.074\,\textdegree{}C over 6.5 hours, returned to setpoint 1.6 minutes after being opened, and was unaffected by a 2.8\,\textdegree{}C rise in room temperature. Six of the ten requirements were met and R2 was partly met. None of the remaining shortfalls came from the controller. They came from where heat enters the chamber, a CO\textsubscript{2} connection that leaked, and loose wiring.

Holding the incubator air at setpoint did not hold the sample at setpoint. With the air at 37\,\textdegree{}C, the water in a well sat at 34.4\,\textdegree{}C, and raising the air temperature warmed every well by the same amount while leaving a 2.5\,\textdegree{}C difference across the plate, so no air temperature could bring every well to 37\,\textdegree{}C. Heating the plate from below solved this, bringing the centre well to within 0.19\,\textdegree{}C of 37\,\textdegree{}C with the air unchanged, even though the base heater ran at a fixed duty without feedback. Putting the base under closed-loop control, with a sensor at the plate, would allow the sample itself to be held at 37\,\textdegree{}C. Moving the base busbars or splitting the base into two zones would then even out the remaining difference across the plate.


The humidity was limited by the same temperature offset. With no humidified supply, the only source of moisture was the water in the wells, and air can only be saturated at the temperature of the water beneath it. The wells sat at 34.5\,\textdegree{}C against 39\,\textdegree{}C air, which caps the chamber at 78\,\%, and 77.3\,\% was measured, short of the 90\,\% of R3. Adding more water does not raise this limit. Reaching 90\,\% needs either a heated water reservoir inside the chamber, held at 37\,\textdegree{}C or warmer, or a continuous humidified air supply separate from the CO\textsubscript{2} valve, since only about a twentieth of the air entering the chamber comes through the CO\textsubscript{2} supply. The external supply is preferred, because standing water inside the chamber cannot be adjusted during a run and is not acceptable for the BSL-3 work this incubator is intended for. In either case the sample and walls must stay close to the air temperature, since at 39\,\textdegree{}C and 90\,\% humidity the dew point is about 37\,\textdegree{}C, and any colder surface, would collect condensation.

The models were built before the enclosure design was final, so their values were wrong: they predicted 2.5 times the measured chamber gain, about twice the glass gain, and a gas volume 2.1 times too large. Their form was right. Both thermal systems behaved as stable, linear, first order systems with negligible delay, so only the gain and time constant needed updating from step tests, and the controller design did not change. With only one build possible, designing in simulation and then measuring the real system was the right approach.

The two temperature loops affect each other strongly. The wall heaters warm the glass 84\,\% as much as they warm the air, which lowers the gain the chamber loop sees from 51 to 32.6 while the glass loop is running. Tuning without this correction would have made the chamber loop 1.6 times more aggressive than intended. With it, the glass held $48.00 \pm 0.10$\,\textdegree{}C while the chamber moved by 3.8\,\textdegree{}C. 

The temperature results are limited by the sensors. The converter resolves 0.034\,\textdegree{}C and the chamber reading took only eight different values over 80 minutes, so the quoted standard deviation reflects the sensor resolution. The PT1000 reference is accurate to about $\pm$0.22\,\textdegree{}C, most of the $\pm$0.3\,\textdegree{}C that R1 allows, and was calibrated at 0\,\textdegree{}C with the thermistors matched at 23.6\,\textdegree{}C, both well below the operating point. Tighter claims need better sensors and a calibration near 37\,\textdegree{}C.

The incubator cost R\,13\,527, an order of magnitude below the commercial systems reviewed, and its design, firmware and test data are documented so that it can be rebuilt or modified. Of the gap identified in Section~\ref{sec:review-summary}, the project delivered compact, modifiable hardware with several control loops on one embedded platform. It did not deliver long-duration performance: the CO\textsubscript{2} loop was never tested with gas, the longest continuous run was 6\,h\,33\,min against the 24 hours of R8, and focal drift was not measured.
